\clearpage
\section{Hydrodynamics}
\begin{colbox}
	\begin{enumerate}
		\item The Maxwell model for viscoelastic fluids corresponds to a spring and dashpot in series. Its rhelogy is given by \begin{equation}
			\ins{1+\tau\del_t}\tilde{\sigma}=\eta\dot{\gamma}
		\end{equation}
		where \(\tau = \eta / G\) is the relaxation time, \(\tilde{\sigma}\) is the shear stress, \(\eta\) is the ciscosity and \(\dot{\gamma}\) is the shear strain rate. Show explicitly that the equation is solved by \begin{equation}
			\tilde{\sigma} = G\int\limits_0^t \dl t' \dot{\gamma}(t') e^{-\ins{t-t'}/\tau}
		\end{equation}
		demonstrate the solid/fuid behaviors for \(t\ll\tau\) and \(t\gg\tau\).
		\item The Kelvin-Voigt model corresponds to a spring and dashpot in parallel. What is the rheology of this model \(\ins{\sigma(\gamma,\dot{\gamma})}\)? How does this material behave at long times?
		\item Consider a rigid body rotation, whose velocity field is described by \(\vec{v} = \vec{\Omega}\times\vec{r}\) with \(\vec{\Omega}\) the angular velocity. Show that the tensor \(\del_\alpha v_\beta\) is antisymmetric in this case.
	\end{enumerate}
\end{colbox}
\subsection{Maxwell model for viscoelastic fluids}
inserting \(\eta = \tau G\) and the given integral expression
\begin{align}\label{eq1}
	\ins{1+\tau\del_t}\int\limits_0^t \dl t' \dot{\gamma}(t') e^{-\ins{t-t'}/\tau} = \tau\dot{\gamma}
\end{align}
pulling a constant w.r.t the integral out
\begin{equation}
	\ins{1+\tau\del_t}e^{-t/\tau}\int\limits_0^t \dl t' \dot{\gamma}(t') e^{t'/\tau} = \tau\dot{\gamma}
\end{equation}
If we designate \(A(t) = e^{-t/\tau}\) and \(B(t) = \int\limits_0^t \dl t' \dot{\gamma}(t') e^{t'/\tau}\) then the equation becomes
\begin{equation}
	\ins{1+\tau\del_t}A(t)B(t) = A(t)B(t) + \tau\ins{\diff{A(t)}{t}B(t)+A(t)\diff{B(t)}{t}}
\end{equation}
and since
\begin{equation}
	\diff{A}{t} = -\frac{1}{\tau}e^{-t/\tau} = -\frac{1}{\tau}A(t)
\end{equation}
we can rewrite
\begin{align}
	\ins{1+\tau\del_t}A(t)B(t) &= A(t)B(t) + \tau\ins{-\frac{1}{\tau}A(t)B(t)+A(t)\diff{B(t)}{t}} \\
	&= A(t)\ins{B(t) + \tau\ins{-\frac{1}{\tau}B(t)+\diff{B(t)}{t}}} \\
	&= \tau A(t)\diff{B(t)}{t}
\end{align}
inserting into equation~\ref{eq1}
\begin{equation}
	e^{-t/\tau}\del_t\int\limits_0^t \dl t' \dot{\gamma}(t') e^{t'/\tau} = \dot{\gamma}
\end{equation}
but
\begin{equation}
	\del_t\int\limits_0^t \dl t' \dot{\gamma}(t') e^{t'/\tau} = \dot{\gamma}(t)e^{t/\tau}
\end{equation}
so eventually
\begin{equation}
	\dot{\gamma} = \dot{\gamma}
\end{equation}
so the expression for \(\tilde{\sigma}\) solves the equation.
\subsubsection{\(t\ll\tau\)}
In this limit
\begin{equation}
	\tilde{\sigma} = G\int\limits_0^t \dl t' \dot{\gamma}(t') e^{-\ins{t-t'}/\tau} \simeq  G\int\limits_0^t \dl t' \dot{\gamma}(t') = G\ins{\gamma(t)-\gamma(0)}
\end{equation}
and since the stress is linear in strain, as we know from the lectures, the system behaves as an ideal elastic solid.
\subsubsection{\(t\gg\tau\)}
In this limit the exponential term in the integral expression for \(\tilde{\sigma}\) decays quickly and earlier values of \(\dot{\gamma}\) are cancelled out by the exponential approaching zero. This means we can assume the main contribution comes from very late times so we can treat \(\dot{\gamma}\) as a constant:
\begin{align}
	\tilde{\sigma} &= G\int\limits_0^t \dl t' \dot{\gamma}(t') e^{-\ins{t-t'}/\tau} \simeq G\dot{\gamma}(t)\int\limits_0^t \dl t' e^{-\ins{t-t'}/\tau} \\
	&= \tau G \dot{\gamma}(t) \\
	&= \eta \dot{\gamma}
\end{align}
which is the equation for an ideal viscous liquid as we know from the lectures.
\subsection{Kelvin-Voigt model}
Since the spring and dashpot are connected in parallel, the strain on both of the elements will be the same, the strain comes from two separate contributions:
\begin{align}
	\gamma = \gamma_{spring} = \gamma_{dashpot} && \sigma = \sigma_{spring} + \sigma_{dashpot}
\end{align}
as we know from the lectures
\begin{align}
	\sigma_{spring} = G\gamma && \sigma_{dashpot} = \eta\dot{\gamma}
\end{align}
which gives us
\begin{equation}
	\sigma = G\gamma + \eta\dot{\gamma}
\end{equation}
which is the Kelvin-Voigt rheology. We can treat this as a separable differential equation in \(\gamma\) whose solution is
\begin{equation}
	\gamma = \frac{\sigma_0}{G} + ce^{-Gt/\eta}
\end{equation}
where \(\sigma_0\) is some stress value. For very large times
\begin{equation}
	\gamma = \frac{\sigma_0}{G}
\end{equation}
\subsection{Antisymmetric tensor}
I'll rewrite the velocity using levi-civita symbols
\begin{equation}
	v_i = \epsilon_{ijk}\Omega_j r_k
\end{equation}
applying the derivative
\begin{equation}
	\del_\alpha v_\beta =  \epsilon_{\beta jk}\Omega_j\ins{\del_\alpha r_k} = \epsilon_{\beta jk} \Omega_j \delta_{\alpha k} = \epsilon_{\beta j \alpha}\Omega_j = -\epsilon_{\alpha j \beta}\Omega_j = -\del_\beta v_\alpha
\end{equation}
therefore it's antisymmetric.